\makesectionpage[figures/Hyperparameter.png]{Hyperparameter}


\begin{frame}{Hyperparameter: Grundlagen}
\begin{columns}
\column{0.55\textwidth}
\begin{itemize}
  \item \textbf{Hyperparameter} beeinflussen die \textbf{Struktur} eines
        Lernmodells.
  \item Sie werden bei der \textbf{Erstellung} des Modell-Objektes
        mitgegeben.
  \item Die \textbf{Parameter} des Modells passt der Lernalgorithmus
        dagegen während des \textbf{Trainings} an.
\end{itemize}

\column{0.40\textwidth}
\begin{blockC1}{Beispiele}
\begin{itemize}
  \item Anzahl der \textbf{Layer} und \textbf{Neuronen} im neuronalen Netz.
  \item \textbf{Kernel}, $\gamma$ und $C$ bei einer SVM.
  \item Maximale \textbf{Tiefe} eines Entscheidungsbaumes.
\end{itemize}
\end{blockC1}
\end{columns}
\source{\cite{seemann2024mlaz}}
\end{frame}


\begin{frame}{Auswahl der besten Hyperparameter}
\begin{columns}
\column{0.35\textwidth}
\centering
\twemoji[width=0.9\linewidth]{bullseye}

\column{0.60\textwidth}
Wie finden wir die \textbf{besten} Hyperparameter?

\vspace{0.4em}
\begin{itemize}
  \item \textbf{Idee:} Wir probieren viele verschiedene Parameter aus.
  \item Wir \textbf{vergleichen} die Ergebnisse.
  \item Wir wählen das Modell mit den \textbf{besten} Ergebnissen.
\end{itemize}
\end{columns}
\source{\cite{seemann2024mlaz}}
\end{frame}

\begin{frame}{Hyperparameter: GridSearchCV -- Grundidee}
\begin{columns}
\column{0.55\textwidth}
\begin{itemize}
  \item sklearn bietet die Klasse \textbf{GridSearchCV} zur Bestimmung der
        besten Hyperparameter.
  \item Man gibt ein \textbf{Gitter} (grid) möglicher Werte je
        Hyperparameter vor.
  \item GridSearchCV probiert \textbf{alle Kombinationen} systematisch aus.
  \item Jede Kombination wird bewertet, die \textbf{beste} wird
        zurückgegeben.
\end{itemize}

\column{0.40\textwidth}
\begin{blockC1}{Gitter-Beispiel}
Für eine SVM etwa:
\begin{itemize}
  \item $C \in \{0{,}1,\; 1,\; 10\}$
  \item $\gamma \in \{0{,}01,\; 0{,}1\}$
\end{itemize}
Das ergibt $3 \times 2 = 6$ Kombinationen.
\end{blockC1}
\end{columns}
\source{\cite{seemann2024mlaz}}
\end{frame}

\begin{frame}{Hyperparameter: GridSearchCV -- Cross-Validation}
\begin{columns}
\column{0.55\textwidth}
Das \textbf{CV} in GridSearchCV steht für \textbf{Cross-Validation}
(Kreuzvalidierung):
\begin{itemize}
  \item Die Trainingsdaten werden in $k$ \textbf{Teile} (Folds) zerlegt.
  \item Das Modell wird $k$-mal trainiert -- jeweils auf $k-1$ Teilen, mit
        einem Teil zur \textbf{Validierung}.
  \item Der Score ist der \textbf{Mittelwert} über alle $k$ Durchläufe.
\end{itemize}

\column{0.40\textwidth}
\begin{blockC1}{Vorteil}
Die Bewertung ist \textbf{robuster}, da nicht von einer einzelnen
zufälligen Aufteilung abhängig.
\end{blockC1}
\end{columns}
\source{\cite{seemann2024mlaz}}
\end{frame}

\begin{frame}[fragile]{Hyperparameter: GridSearchCV in Python}
\begin{columns}
\column{0.60\textwidth}
\begin{minted}[fontsize=\footnotesize]{python}
from sklearn.model_selection import GridSearchCV
from sklearn.svm import SVC

param_grid = {
    'C': [0.1, 1, 10],
    'gamma': [0.01, 0.1],
    'kernel': ['rbf']
}
grid = GridSearchCV(SVC(), param_grid, cv=5)
grid.fit(X_train, y_train)
print(grid.best_params_)
print(grid.best_score_)
\end{minted}

\column{0.35\textwidth}
\begin{blockC1}{Wichtige Attribute}
\begin{itemize}
  \item \texttt{param\_grid}: das \textbf{Gitter} der Werte.
  \item \texttt{cv}: Anzahl der \textbf{Folds}.
  \item \texttt{best\_params\_}: beste \textbf{Kombination}.
  \item \texttt{best\_score\_}: zugehöriger \textbf{Score}.
\end{itemize}
\end{blockC1}
\end{columns}
\source{\cite{seemann2024mlaz}}
\end{frame}

\begin{frame}[fragile]{GridSearchCV: Bearbeitungs-Pipeline}
\begin{columns}
\column{0.60\textwidth}
Eine \textbf{Pipeline} legt fest, welche Schritte nacheinander auf den
Daten ausgeführt werden; hier erst der \texttt{StandardScaler}, dann der
\texttt{RandomForestClassifier}. Jede Komponente bekommt einen
\textbf{Namen}:
\begin{minted}[fontsize=\footnotesize]{python}
from sklearn.pipeline import Pipeline
from sklearn.preprocessing import StandardScaler
from sklearn.ensemble import RandomForestClassifier

pipeline = Pipeline([
    ("scaler", StandardScaler()),
    ("rf", RandomForestClassifier())
])
\end{minted}

\column{0.35\textwidth}
\begin{blockC1}{Aufbau}
\begin{itemize}
  \item Jede Stufe ist ein \textbf{Paar}: Name und Objekt.
  \item \texttt{scaler}: \textbf{Skalierung} der Merkmale.
  \item \texttt{rf}: der \textbf{Random Forest}.
  \item Die Namen brauchen wir später im \textbf{Parametergitter}.
\end{itemize}
\end{blockC1}
\end{columns}
\source{GridsearchRandomForest.ipynb}
\end{frame}

\begin{frame}[fragile]{GridSearchCV: Parametersuche mit Pipeline}
\begin{columns}
\column{0.60\textwidth}
\textbf{GridSearchCV} probiert verschiedene Konfigurationen aus und gibt die
Kombination mit dem besten Score zurück, automatisch mit
\textbf{k-facher Kreuzvalidierung}:
\begin{minted}[fontsize=\footnotesize]{python}
from sklearn.model_selection import GridSearchCV

clf = GridSearchCV(pipeline, param_grid={
    "rf__n_estimators": range(470,500,1)
})
clf.fit(X, y)
print(clf.best_params_)
\end{minted}

\begin{minted}[frame=none]{text}
{'rf__n_estimators': 472}
\end{minted}

\column{0.35\textwidth}
\begin{blockC1}{Namenskonvention}
Der Schlüssel im Gitter setzt sich zusammen aus:
\[
  \underbrace{\texttt{rf}}_{\text{Komponente}}
  \;\texttt{\_\_}\;
  \underbrace{\texttt{n\_estimators}}_{\text{Parameter}}
\]
Dahinter steht das \textbf{Feld} aller zu testenden Werte.
\end{blockC1}
\end{columns}
\source{GridsearchRandomForest.ipynb}
\end{frame}

\begin{frame}[fragile]{GridSearchCV: Pipeline mit MLP}
\begin{columns}
\column{0.60\textwidth}
Die Pipeline kombiniert \texttt{StandardScaler} und
\texttt{MLPClassifier}. Beim Anlegen lassen sich schon \textbf{einzelne
Parameter} setzen:
\begin{minted}[fontsize=\footnotesize]{python}
from sklearn.pipeline import Pipeline
from sklearn.preprocessing import StandardScaler
from sklearn.neural_network import MLPClassifier

pipeline = Pipeline([
    ("scaler", StandardScaler()),
    ("mlp", MLPClassifier(solver='lbfgs',
                          max_iter=20000))
])
\end{minted}

\column{0.35\textwidth}
\begin{blockC1}{Aufbau}
\begin{itemize}
  \item \texttt{scaler}: \textbf{Skalierung} der Merkmale.
  \item \texttt{mlp}: das \textbf{neuronale Netz}.
  \item Feste Parameter (\texttt{solver}, \texttt{max\_iter}) setzt man
        direkt im \textbf{Konstruktor}.
\end{itemize}
\end{blockC1}
\end{columns}
\source{GridsearchNeuronalesNetz.ipynb}
\end{frame}

\begin{frame}[fragile]{GridSearchCV: Mehrere Parameter optimieren}
  \footnotesize
\begin{columns}
\column{0.60\textwidth}
Mit GridSearchCV lassen sich \textbf{beliebig viele} Parameter optimieren --
alle \textbf{Kombinationen} werden ausprobiert:
\begin{minted}[fontsize=\footnotesize]{python}
from sklearn.model_selection import GridSearchCV

clf = GridSearchCV(pipeline, param_grid={
    "mlp__hidden_layer_sizes":
        [(50,50,50), (30,20,10), (500,10),
         (20,20,20,20), (30,20)],
    "mlp__alpha":
        [10, 5, 3, 2, 1, 0.1, 0.01, 0.001]
})
clf.fit(X, y)
print(clf.best_params_)
\end{minted}

\begin{minted}[frame=none]{text}
{'mlp__alpha': 3,
 'mlp__hidden_layer_sizes': (20, 20, 20, 20)}
\end{minted}

\column{0.35\textwidth}
\begin{blockC2}{\footnotesize Achtung}
Viele Parameter lassen die Zahl der \textbf{Kombinationen} -- und damit die
\textbf{Laufzeit} -- schnell \textbf{explodieren}.
\end{blockC2}
\end{columns}
\source{GridsearchNeuronalesNetz.ipynb}
\end{frame}

\begin{frame}{Classification Challenge: MNIST-Datensatz}
\small
\begin{columns}
\column{0.45\textwidth}
\begin{itemize}
  \item \textbf{MNIST} ist einer der Standarddatensätze zum
        \textbf{Benchmarking} von Klassifikatoren.
  \item Er umfasst \textbf{70.000} handgeschriebene Ziffern samt
        \textbf{Label}.
  \item Feste Aufteilung: \textbf{60.000} Trainings- und \textbf{10.000}
        Testdatensätze.
  \item Die feste Teilung schließt \textbf{zufallsbasierte} Schwankungen in
        den Testergebnissen aus.
\end{itemize}

\column{0.55\textwidth}
\centering
\includegraphics[width=\linewidth]{figures/mnistexamples.png}
\end{columns}
\source{Bildquelle: \cite{steppan_mnist_examples}}
\end{frame}

\begin{frame}{Classification Challenge: MNIST-Datensatz}
\begin{columns}
\column{0.55\textwidth}
\begin{itemize}
  \item Jedes Zeichen ist ein Bild mit \textbf{28$\times$28 Pixeln}.
  \item Jedes Pixel hat einen Wert zwischen \textbf{0 und 255}.
  \item Insgesamt \textbf{70.000} solcher Bilder, als gezippte Dateien.
  \item Entpackt haben die Daten eine Größe von \textbf{52 MB}.
\end{itemize}

\column{0.40\textwidth}
\begin{blockC1}{Erreichbare Trefferquote}
\begin{itemize}
  \item In Fachartikeln: über \textbf{99\%} mit sehr speziellen Verfahren
        und aufwändiger Datenaufbereitung.
  \item Mit unseren aktuellen Klassifikatoren realistisch: etwa \textbf{98\%}.
\end{itemize}
\end{blockC1}
\end{columns}
\source{\cite{seemann2024mlaz}}
\end{frame}

\begin{frame}{Übung: Ziffernerkennung mit GridSearch}
\begin{columns}
\column{0.60\textwidth}
Verwenden Sie die Vorlage \textbf{ProjektZiffernerkennung.ipynb}. Sie
enthält bereits das Einlesen und Aufbereiten der Trainings- und Testdaten.

\vspace{0.4em}
\begin{itemize}
  \item Erweitern Sie das Notebook, sodass statt eines einfachen Trainings
        eine \textbf{GridSearch} durchgeführt wird.
  \item Erstellen Sie verschiedene Notebooks für die verschiedenen
        \textbf{Klassifikationsverfahren} und suchen Sie jeweils die passenden
        Parameter.
\end{itemize}

\vspace{0.4em}
\textbf{Frage:} Welche beste \textbf{Prognosegüte} erreichen Sie auf allen
\textbf{10.000} Testdatensätzen?

\column{0.35\textwidth}
\centering
\twemoji[width=0.9\linewidth]{laptop}
\end{columns}
\source{Vorlage: ProjektZiffernerkennung.ipynb}
\end{frame}

\makequizpage{\QUIZURL}{\QUIZQRCODE}{\QUIZIMAGE}
\makequestionspage{\FRAGENIMAGE}